%!TEX TS-program = PGF
%!TEX encoding = UTF-8
\documentclass{standalone}
\input{tikzstandalone}
\begin{document}

\def\courbepHKe{1e-14}% produit ionique de l'eau
\def\courbepHKaa{10^(-6.3)}% Ka1
\def\courbepHKab{10^(-10.3)}% Ka2
\def\courbepHCa{0.10} %becher
\def\courbepHVa{10.0} % volume d'acide (mL)
\def\courbepHCb{0.10} % burette
\pgfmathsetmacro{\Veq}{\courbepHCb*\courbepHVa/\courbepHCb}
\pgfmathsetmacro{\courbepHka}{\courbepHKaa}
\pgfmathsetmacro{\courbepHkb}{\courbepHKaa*\courbepHKab}

\begin{tikzpicture}[
	declare function={%déclaration de fonctions :
	Oxo(\pH)=10^(-\pH);% h concentration des ions H3O+
	x0(\pH)=1/(1+\courbepHka/Oxo(\pH)+\courbepHkb/(Oxo(\pH)^2));%fraction acide
	x1(\pH)=x0(\pH)*\courbepHka/Oxo(\pH);% fraction base
	x2(\pH)=x0(\pH)*\courbepHkb/(Oxo(\pH)^2);% fraction base
	Y(\pH)=x1(\pH) + 2*x2(\pH);% moment
	w(\pH)=\courbepHKe/Oxo(\pH)-Oxo(\pH); % w-h
	frac_base(\pH)=Y(\pH) + w(\pH)/\courbepHCa;
	V_verse(\pH)=frac_base(\pH)*\Veq;% volume de base versée
}]
	
	\begin{axis}[% mise en forme du graphique :
		xmin=0,xmax=14,
		ymin=0,ymax=1,
		xlabel={volume $V_b$ versé (\si{mL})},
		ylabel={\pH},
		ytick distance=2,
		minor y tick num=1,% nombre de graduations secondaires intermédiaires
		minor x tick num=4,% nombre de graduations secondaires intermédiaires
		grid,
		grid=both,% grilles primaire et secondaire
		title={titrage pH-métrique d'un di-acide faible},
		axis line shift=3pt,
		xtick align=outside,
		ytick align=outside,
		tickpos=left,
	]
		\addplot[domain=0:14, thick, samples=200] ({x},{x0(x)});% volume versé en x, et pH en y
		\addplot[domain=0:14, thick, samples=200] ({x},{x1(x)});% volume versé en x, et pH en y
		\addplot[domain=0:14, thick, samples=200] ({x},{x2(x)});% volume versé en x, et pH en y
	\end{axis}
	
\end{tikzpicture}

\end{document}
